L'auto-évaluation reprend la grille institutionnelle C1--C9.

\begin{longtable}{p{0.15\textwidth}p{0.37\textwidth}p{0.37\textwidth}}
\caption{Auto-évaluation des compétences mobilisées pendant le stage.}
\label{tab:cs_self_assessment}\\
\hline
Compétence & Points forts & Points d'amélioration \\
\hline
\endfirsthead
\hline
Compétence & Points forts & Points d'amélioration \\
\hline
\endhead
C1 --- Complexité &
Décomposition traçable du problème : DFT $\rightarrow$ Hamiltonien effectif
$\rightarrow$ action euclidienne $\rightarrow$ SCHA $\rightarrow$
susceptibilité. Hypothèses, échelles, incertitudes, sensibilité au volume et
contrôles de convergence sont explicités. &
Comparer systématiquement la fermeture aux dynamiques moléculaires ou
Monte-Carlo et quantifier l'incertitude de modèle, notamment près de la
transition. \\
\hline
C2 --- Métier ingénieur &
Approfondissement autonome de la physique des phonons, des Hamiltoniens
effectifs, de la théorie des champs à température finie et des méthodes
numériques. Construction d'une chaîne de calcul reproductible. &
Consolider l'expertise en calcul de matériaux à température finie et
implémenter une correction post-SCHA auto-cohérente. \\
\hline
C3 --- Innover, entreprendre &
Reformulation du problème autour de diagnostics falsifiables : séparation
volume/masse, comparaison Mayer, test post-SCHA et prémices VCA pour BST. &
Transformer ces pistes en plan de développement validé par des calculs de
référence et définir un périmètre de diffusion du solveur. \\
\hline
C4 --- Création de valeur &
Production d'un intermédiaire analytique et numérique réutilisable pour le
SPMS, avec scripts, contrôles, paramètres et documentation destinés au dépôt
GitHub. &
Définir avec les utilisateurs du laboratoire des indicateurs d'usage du dépôt
et prioriser les extensions selon leur coût et leur apport expérimental. \\
\hline
C5 --- Interculturel &
Échanges scientifiques avec un superviseur anglophone, lecture critique
d'articles en anglais, rédaction d'un abstract anglais et usage de références
internationales dans un contexte de recherche. &
Compétence peu mobilisée directement : développer une collaboration
internationale ou une mobilité et expliciter davantage les enjeux
interculturels. \\
\hline
C6 --- Numérique &
Développement et audit de solveurs Python, automatisation de la validation,
traçabilité Git, contrôle des sorties et déclaration de l'assistance au codage.
& \\
\hline
C7 --- Convaincre &
Rapport argumenté, figures comparatives, distinction explicite entre prédiction
et calibration, préparation d'une soutenance adaptée à un public non spécialiste
et échanges réguliers avec l'encadrement. &
Élargir la confrontation des résultats et du support à des interlocuteurs
expérimentaux ou industriels afin d'éprouver la transférabilité du message et
les usages concrets du solveur. \\
\hline
C8 --- Équipe, projet &
Organisation du travail par jalons, échanges avec les encadrants et membres du
SPMS, intégration de retours scientifiques dans un livrable commun et
traçable. &
Le stage ne comportait pas de responsabilité hiérarchique directe : développer
la coordination d'un sous-projet collectif et la formalisation partagée des
décisions. \\
\hline
C9 --- Éthique, soutenabilité &
Intégrité scientifique : délimitation claire de ce qui a été fait et des
contributions personnelles ; sources, données reconstruites, extrapolations et
calibrations sont distinguées. Les limites de modèle, les incertitudes et la
sobriété numérique sont documentées. Le travail sur BST s'inscrit aussi dans la
recherche de matériaux ferroélectriques sans plomb, dans la perspective de
réduire le recours au Pb des PZT. & \\
\hline
\end{longtable}
